\section{Illustrative discordant cases}
\label{sec:heldout-cases}

The following four examples were selected after generation by a fixed deterministic rule: for each label contrast, choose the lowest numeric task ID in each direction in which both programs were observed and eligibility-gated, retaining one role-labelled--neutral discordance in each direction. They are descriptive examples outside the four primary tests. They do not estimate discordance frequencies or identify a hidden reasoning mechanism.

For \texttt{BigCodeBench/428}, the role-labelled single-call program passed while the neutral program failed. Both programs performed the requested outer merge and scaling of numeric features from \texttt{df1}. The role-labelled program scaled a copy of \texttt{df1} before one merge and preserved the left-dataframe columns first. The neutral program merged first, then dropped and reinserted the scaled columns in a second merge, producing the order \texttt{id, feature4, feature5, feature1, feature2, feature3}. The evaluator's \texttt{test\_case\_1} required \texttt{id, feature1, feature2, feature3, feature4, feature5} and failed with that list mismatch. The task prose specifies the operations but does not explicitly specify column order; the test supplies a stricter output contract.

For \texttt{BigCodeBench/71}, the neutral single-call program passed while the role-labelled program failed. The role-labelled implementation used sample standard deviation, \texttt{Series.std(ddof=1)}, while the neutral implementation and canonical solution used population standard deviation, \texttt{np.std} with \texttt{ddof=0}. The role-labelled program consequently failed \texttt{test\_case\_2}, \texttt{test\_case\_3}, \texttt{test\_case\_4}, and \texttt{test\_case\_5}; the recorded discrepancies were respectively $11.900380145988935$ versus $11.017611134090766$, $8.568005268772557$ versus $8.01463505095522$, NaN versus $0.0$, and $47.95684491664328$ versus $46.07544623291379$. “Standard deviation” is underspecified in the prose, while the executable contract fixes the population convention.

For \texttt{BigCodeBench/167}, the role-labelled three-call program passed while the neutral program failed. The neutral implementation selected \texttt{max(2,min(5,num\_types))} stacked segments. It therefore produced 50 patches for \texttt{num\_types=10} and 2 for \texttt{num\_types=1}, whereas the role-labelled implementation used exactly \texttt{num\_types} segments and produced 100 and 1. The native failures were \texttt{test\_case\_3}, $50\mathrel{\ne}100$, and \texttt{test\_case\_4}, $2\mathrel{\ne}1$. The prose describes categories and segmented values but leaves the segment count open; the tests require it to equal \texttt{num\_types}.

For \texttt{BigCodeBench/31}, the neutral three-call program passed while the role-labelled program failed. The neutral implementation preserved the first-occurrence order of the \texttt{Counter}; the role-labelled implementation sorted by descending frequency and then lexicographically. On the test text containing \texttt{\$is} three times, \texttt{\$second} once, and \texttt{\$sentence} twice, the evaluator expected values $3,1,2$ in first-occurrence category order. Sorting moved \texttt{\$sentence} ahead of \texttt{\$second}, and \texttt{test\_case\_2} reported ``Expected value '1.0', but got '2'.'' The prose requires a frequency plot without specifying category order, so this is an executable-order mismatch rather than evidence of a causal role effect.

All descriptions use the recorded prompts, programs, and native assertions. They describe observed code behavior only; they do not infer private reasoning, latent agents, or representative population frequencies.
